% Overriding the theorem counter for this chapter to match the 14.1, 14.2 convention
\renewcommand{\thetheorem}{\thechapter.\arabic{theorem}}
\setcounter{chapter}{13} % Sets chapter counter so the current chapter is 14
\setcounter{theorem}{0}

\chapter{Sums of Vector Spaces}

% def:14.1
\begin{definition}[Sum of Subspaces]
\label{def:sum_of_subspaces}
Let $V$ be a vector space over $K$, and let $U, W \subseteq V$ be linear subspaces. We define the \qt{sum} of $U$ and $W$, denoted by $U + W$, as the set:
\[
U + W := \{u + w \mid u \in U, w \in W\}.
\]
In a similar way, we can define sums of more than two subspaces. If $U_1, \dots, U_r \subseteq V$ are subspaces, we define their total sum as:
\[
U_1 + \dots + U_r := \{u_1 + \dots + u_r \mid u_1 \in U_1, \dots, u_r \in U_r\}.
\]
\end{definition}

% prop:14.2
\begin{proposition}[Properties of Sums]
\label{prop:properties_of_sums}
Let $U, W \subseteq V$ be subspaces. Then, the following properties hold:
\begin{enumerate}
    \item $U + W = \operatorname{Sp}(U \cup W)$. In particular, $U + W \subseteq V$ is a subspace, and $U$ and $W$ are themselves subspaces of $U + W$.
    \item If $U$ and $W$ are finite-dimensional, then so is $U + W$. Moreover, one can systematically find a basis for $U + W$ by following these sequential steps: first, choose a basis $(v_1, \dots, v_k)$ for $U \cap W$ (where $k = \dim(U \cap W)$). Second, extend this to a basis $(v_1, \dots, v_k, u_1, \dots, u_{l-k})$ for $U$ (where $l = \dim U$). Third, extend the initial intersection basis to a basis $(v_1, \dots, v_k, w_1, \dots, w_{m-k})$ for $W$ (where $m = \dim W$). Finally, the combined elements 
    \[(v_1, \dots, v_k, u_1, \dots, u_{\ell-k}, w_1, \dots, w_{m-k})\] 
    form a basis for $U + W$.
\end{enumerate}
\end{proposition}

\begin{proof}
We will explicitly prove the second statement. By definition, the subspace $U + W$ is generated by the combined vectors $(v_1, \dots, v_k, u_1, \dots, u_{l-k}, w_1, \dots, w_{m-k})$. It remains to show that these vectors are linearly independent.

Suppose there exists a linear combination equal to the zero vector:
\begin{equation}
\label{eq:linear_comb_sum}
\alpha_1 v_1 + \dots + \alpha_k v_k + \beta_1 u_1 + \dots + \beta_{\ell-k} u_{\ell-k} + \gamma_1 w_1 + \dots + \gamma_{m-k} w_{m-k} = 0.
\end{equation}
We can rearrange this equation to isolate the terms belonging to $U$ and $W$:
\[
\underbrace{\beta_1 u_1 + \dots + \beta_{\ell-k} u_{\ell-k}}_{\in U} = \underbrace{-\alpha_1 v_1 - \dots - \alpha_k v_k - \gamma_1 w_1 - \dots - \gamma_{m-k} w_{m-k}}_{\in W}.
\]
Let us denote this common vector by $x$. Since the left side is entirely in $U$ and the right side is entirely in $W$, we have $x \in U$ and $x \in W$, which implies $x \in U \cap W$. 

Because $(v_1, \dots, v_k)$ is a basis for $U \cap W$, there must exist scalars $\delta_1, \dots, \delta_k \in K$ such that:
\[
x = \delta_1 v_1 + \dots + \delta_k v_k.
\]
Equating this with our original expression for $x$ in $U$, we obtain:
\[
\beta_1 u_1 + \dots + \beta_{\ell-k} u_{\ell-k} = \delta_1 v_1 + \dots + \delta_k v_k \implies \sum_{i=1}^{\ell-k} \beta_i u_i - \sum_{j=1}^k \delta_j v_j = 0.
\]
However, the vectors $(v_1, \dots, v_k, u_1, \dots, u_{\ell-k})$ form a basis for $U$, meaning they are strictly linearly independent. Consequently, all coefficients in this sum must be zero. In particular, $\beta_1 = \dots = \beta_{\ell-k} = 0$.

Substituting these zero values back into our initial equation \eqref{eq:linear_comb_sum}, the terms involving $u_i$ vanish, leaving:
\[
\alpha_1 v_1 + \dots + \alpha_k v_k + \gamma_1 w_1 + \dots + \gamma_{m-k} w_{m-k} = 0.
\]
Since the vectors $(v_1, \dots, v_k, w_1, \dots, w_{m-k})$ form a basis for $W$, they are linearly independent. This immediately implies that $\alpha_1 = \dots = \alpha_k = 0$ and $\gamma_1 = \dots = \gamma_{m-k} = 0$. 

Because all coefficients $\alpha_i$, $\beta_j$, and $\gamma_k$ have been forced to zero, the combined family of vectors is linearly independent, and therefore constitutes a valid basis for $U + W$.
\end{proof}

% thm:14.3
\begin{theorem}[Dimension Formula for Sums]
\label{thm:dimension_formula_sums}
Let $V$ be a vector space, and let $U, W \subseteq V$ be finite-dimensional subspaces. Then, the following dimension formula holds:
\[
\dim(U + W) = \dim U + \dim W - \dim(U \cap W).
\]
\end{theorem}

\begin{proof}
Using the basis constructed in Proposition \ref{prop:properties_of_sums}, we can directly count the number of basis vectors for $U + W$:
\begin{align*}
\dim(U + W) &= k + (\ell - k) + (m - k) \\
&= \ell + m - k \\
&= \dim U + \dim W - \dim(U \cap W).
\end{align*}
\end{proof}

% cor:unnumbered
\begin{corollary}[Equivalent Conditions for Trivial Intersection]
\label{cor:trivial_intersection}
Let $V$ be a vector space, and let $U, W \subseteq V$ be subspaces. The following statements are equivalent:
\begin{enumerate}
    \item $U \cap W = \{0\}$.
    \item For every $v \in U + W$, there exist unique vectors $u \in U$ and $w \in W$ such that $v = u + w$.
    \item Every $x \in U \cap W$ can be uniquely written as $x = u + w$ with $u \in U$ and $w \in W$.
    \item The equality $0 = u + w$ with $u \in U$ and $w \in W$ implies that $u = 0$ and $w = 0$.
\end{enumerate}
\end{corollary}

\begin{proof}
We prove the main equivalences by showing the cycle of implications \textbf{(a)} $\implies$ \textbf{(b)} $\implies$ \textbf{(d)} $\implies$ \textbf{(a)}. (The equivalence of \textbf{(c)} follows naturally from the same uniqueness constraints).

\textbf{(a) $\implies$ (b):} Suppose that $v \in U + W$ can be written in two ways: $v = u + w = u' + w'$, where $u, u' \in U$ and $w, w' \in W$. Rearranging this equation yields $u - u' = w' - w$. The left side is in $U$ and the right side is in $W$, which means the vector $u - u'$ resides in the intersection $U \cap W$. Since we assume $U \cap W = \{0\}$, it follows that $u - u' = 0$ and $w' - w = 0$. Therefore, $u = u'$ and $w = w'$, proving uniqueness.

\textbf{(b) $\implies$ (d):} Consider the zero vector, which can obviously be written as $0 = 0_U + 0_W$. If we also have $0 = u + w$ for some $u \in U$ and $w \in W$, the assumed uniqueness in statement \textbf{(b)} immediately forces $u = 0$ and $w = 0$.

\textbf{(d) $\implies$ (a):} Let $x \in U \cap W$. We can express the zero vector as $0 = x + (-x)$. Since $x \in U \cap W$, we know $x \in U$ and $-x \in W$. If statement \textbf{(d)} holds, this equality strictly implies $x = 0$. Consequently, the only element in the intersection is the zero vector, so $U \cap W = \{0\}$.
\end{proof}

\setcounter{theorem}{3} % Resets the theorem counter so the next environment is labeled 14.4

% def:14.4
\begin{definition}[Complement of a Subspace]
\label{def:complement}
Let $V$ be a vector space over $K$, and let $U \subseteq V$ be a subspace. A subspace $W \subseteq V$ is called a \qt{complement} of $U$ if it satisfies the following two conditions:
\[
U + W = V \quad \text{and} \quad U \cap W = \{0\}.
\]
\end{definition}

\begin{remark}
Note that if $W$ is a complement of $U$, then $U$ and $W$ satisfy all of the five equivalent statements derived from Corollary \ref{cor:trivial_intersection}.
\end{remark}

% prop:14.5
\begin{proposition}[Existence of a Complement]
\label{prop:existence_complement}
Let $V$ be a finite-dimensional vector space, and let $U \subseteq V$ be a subspace. Then, there exists a subspace $W \subseteq V$ which is a complement to $U$.
\end{proposition}

\begin{proof}
Choose a basis $\mathcal{B}_U = (u_1, \dots, u_{\ell})$ for $U$, where $\ell = \dim U$. By the Steinitz Exchange Lemma, we can extend this to a full basis for $V$:
\[
\mathcal{B}_V = (u_1, \dots, u_{\ell}, w_1, \dots, w_m),
\]
where $\ell + m = \dim V$. We define $W := \operatorname{Sp}(w_1, \dots, w_m)$. 

We claim that $W$ is a valid complement of $U$. Indeed, since the combined basis spans the entirety of $V$, it is clear that $U + W = V$. To see that $U \cap W = \{0\}$, let $x \in U \cap W$. Because $x \in U$, we can write $x = \sum_{i=1}^l a_i u_i$. Because $x \in W$, we can also write $x = \sum_{j=1}^m b_j w_j$. Equating these expressions yields:
\[
\sum_{i=1}^{\ell} a_i u_i = \sum_{j=1}^m b_j w_j \implies \sum_{i=1}^{\ell} a_i u_i - \sum_{j=1}^m b_j w_j = 0.
\]
Since $\mathcal{B}_V$ is a basis for $V$, all of its constituent vectors are strictly linearly independent. This forces all scalar coefficients to be zero: $a_i = 0$ for all $i$ and $b_j = 0$ for all $j$. Thus, $x = 0$, proving that $U \cap W = \{0\}$.
\end{proof}